\documentclass[11pt]{article}
\usepackage[utf8]{inputenc}
\usepackage[T1]{fontenc}
\usepackage{lmodern}
\usepackage{amsmath,amssymb}
\usepackage{graphicx}
\usepackage{booktabs}
\usepackage{ulem}
\usepackage{fixltx2e}
\usepackage[margin=20mm,a4paper]{geometry}
\usepackage{hyperref}
\hypersetup{colorlinks=true,linkcolor=blue,urlcolor=blue}
\usepackage{listings}
\lstset{basicstyle=\ttfamily\small,breaklines=true}
\usepackage{kotex}
\title{My Document Converter 샘플}

\date{}

\begin{document}
\maketitle

\section*{소개 (Introduction)}\label{소개-introduction}

Vimwiki 는 \textbf{굵게}, \emph{기울임}, \verb|코드| 와
\href{https://example.com}{링크} 를 지원합니다.
This line is in English.

\subsection*{목록 (Lists)}\label{목록-lists}

\begin{itemize}
\item 첫 번째 항목
\item 두 번째 항목
\begin{itemize}
\item 중첩된 항목
\end{itemize}
\item 세 번째 항목
\end{itemize}

\begin{enumerate}
\item 순서가 있는 항목
\item 두 번째
\item 세 번째
\end{enumerate}

\begin{itemize}
\item $\boxtimes$ 끝난 일
\item $\square$ 남은 일
\end{itemize}

\subsection*{코드 (Code)}\label{코드-code}

\begin{lstlisting}[language=js]
function greet(name) {
  console.log('Hello, ' + name)
}
\end{lstlisting}

\subsection*{표 (Table)}\label{표-table}

\begin{table}[htbp]
\centering
\begin{tabular}{ll}
\toprule
형식 & 확장자 \\
\midrule
Markdown & .md \\
HTML & .html \\
DOCX & .docx \\
\bottomrule
\end{tabular}
\end{table}

\end{document}
